\section{Supplementary Methodology and Detailed Results}
\label{app:detail}

\subsection{Bottom-up Validation of the Residential Building Stock}
\label{sec:method:bottomup}

The top-down heat-demand baseline (Section~\ref{sec:method:data}) rests on the
Hotmaps regional dataset combined with census building-type shares. Because that
baseline is itself a modelled product, we constructed an independent
\emph{bottom-up} estimate of residential heat demand and use it as a validation
cross-check. The bottom-up build is a single pipeline parameterised per country, so that adding a country is a data exercise.

\paragraph{Method.} For each country, every building footprint in the EUBUCCO v0.2 open building dataset \citep{Eubucco2023v02} is classified. The published description of the database covers release v0.1 \citep{MilojevicDupont2023Eubucco}; we use v0.2, which is what extends coverage to the United Kingdom, one of the 29 markets. Each footprint is classified into
one of four types, namely single-family house (SFH), low-rise and high-rise
multi-family (low-rise and high-rise), and non-residential, using sequential rules on
footprint area, floor count and the native EUBUCCO type string. Heated floor area
is the footprint polygon area multiplied by the EUBUCCO floor count and a 0.85
useable-area fraction. Each building is then assigned a space-heating intensity
from the TABULA European building typology \citep{Loga2016Tabula}, indexed by
building class and construction-vintage cohort and multiplied by a
heating-degree-day climate correction and a national retrofit-state blend, with a
domestic-hot-water term added. The resulting per-building annual demand is summed
to a national residential total and reconciled against the Hotmaps baseline.
Countries without a national TABULA dataset use a climate-corrected neighbour as a
proxy (Luxembourg via Belgium \citep{Cyx2011TabulaBelgium}; Finland via Sweden),
with the correction taken from the ratio of Eurostat heating-degree days
\citep{EurostatHDD2025}.

\paragraph{Coverage.} Across the 29 country builds, non-residential buildings
dominate the raw building \emph{count}, between roughly half and two-thirds of all
footprints (garages, saunas and agricultural outbuildings in the open data);
because they carry zero heated floor area, they are excluded from the residential
total. Single-family homes carry 58\% of residential useful heat and multi-family
buildings 42\%, with the single-family share ranging from $\approx$80\% in Ireland
to under 10\% in Malta. The per-vintage intensity gradient that TABULA encodes is
steep. Across countries, the median space-heating intensity falls by roughly a
factor of four from the pre-1945 cohort to post-2020 construction, reflecting
successive thermal building regulations.

\paragraph{Validation result.} Table~\ref{tab:bottomup} and Figure~\ref{fig:bottomup}
compare the bottom-up estimate with the Hotmaps baseline for all 29 countries. On the
raw 2025 snapshot, 19 of the 29 countries fall within the $\pm 15\,\%$ `consistent' tier
and 28 within the $\pm 25\,\%$ acceptance band (only Malta, the smallest stock, lies
outside), while the EU-27+UK+CH aggregate deviates by just $-0.8\,\%$. Three countries
illustrate the spread. Luxembourg ($-20.2\,\%$) sits near the lower edge of the band,
Finland ($-11.8\,\%$) and France ($+11.4\,\%$) inside the tighter tier on opposite sides
of the benchmark. The bottom-up build is reported without any calibration to the
benchmark. This level of agreement, achieved from entirely independent data (building
footprints and a physical typology, set against the Hotmaps top-down accounting),
supports the use of the baseline for the scenario analysis.

A count discards the size of a miss and an aggregate hides it, because the country
deviations carry opposite signs and cancel. On the raw basis the mean absolute deviation
across the 29 countries is $11.6\,\%$, the demand-weighted mean $9.8\,\%$, the
root-mean-square deviation $14.4$ percentage points and the median absolute deviation
$8.5\,\%$, the largest single miss being $35.6\,\%$ at Malta. On the vintage-matched
basis the same four read $12.4$, $10.7$, $15.2$ and $9.6\,\%$. The $-0.8\,\%$ aggregate
is therefore below a tenth of the median country deviation, which is the quantitative
form of the reading we take throughout: the study-area aggregate is the supported
figure, and the country-level verdicts are indicative.

The table's final column reports the vintage-matched deviation, which backcasts each
2025 snapshot to 2015 along the country's historical demand trend before comparison. It
shifts every country downward (the EU-27+UK+CH aggregate moves to $-8.3\,\%$) and is
reported so that the model and the 2015-vintage benchmark are compared on a consistent
year. The remaining gaps are discussed in the Limitations section. In particular,
EUBUCCO construction-year coverage varies sharply between countries (France
approximately 72\,\%, Luxembourg and Finland approximately 0 to 2\,\%), so for
low-coverage countries the result is driven by a stock-weighted fallback intensity in
place of a true per-vintage calculation. Because the aggregate band is tight but the
country tails are not (Luxembourg reaches about $-32\,\%$ on the vintage-matched basis,
and Bulgaria and Romania about $+21$ to $+23\,\%$ on the raw basis), the country-level
verdicts for those tail countries carry materially more uncertainty than the EU-27+UK+CH
aggregate, and should be read as indicative.

% Full 29-country bottom-up validation table, generated by
% code/scripts/make_validation_table.py (carries \label{tab:bottomup}).
% Auto-generated by code/scripts/make_validation_table.py -- do not edit by hand.
% Regenerate: cd code && PYTHONUTF8=1 PYTHONPATH=. python -m scripts.make_validation_table
\begin{table}[H]
\centering
\footnotesize
\setlength{\tabcolsep}{2pt}
\caption{Bottom-up heat demand vs Hotmaps 2015, all 29 countries.}
\label{tab:bottomup}
\begin{tabularx}{\textwidth}{@{}X r S[table-format=4.1] S[table-format=4.1] S[table-format=+3.1] S[table-format=+3.1]@{}}
\toprule
Country & {EUBUCCO} & {Bottom-up} & {Hotmaps 2015} & {Raw} & {Vintage-matched} \\
 & {buildings} & {(TWh\,yr$^{-1}$)} & {(TWh\,yr$^{-1}$)} & {deviation (\%)} & {deviation (\%)} \\
\midrule
United Kingdom & 24{,}861{,}926 & 368.9 & 467.7 & -21.1 & -27.9 \\
Luxembourg & 186{,}171 & 6.6 & 8.3 & -20.2 & -32.3 \\
Estonia & 1{,}141{,}841 & 10.8 & 13.3 & -18.7 & -26.1 \\
Netherlands & 11{,}562{,}667 & 130.7 & 160.7 & -18.6 & -25.6 \\
Austria & 4{,}876{,}313 & 72.1 & 82.4 & -12.4 & -22.3 \\
Finland & 6{,}633{,}364 & 68.9 & 78.1 & -11.8 & -20.9 \\
Czech Republic & 5{,}764{,}786 & 77.9 & 84.6 & -8.0 & -17.9 \\
Denmark & 6{,}152{,}463 & 51.7 & 55.9 & -7.6 & -15.0 \\
Hungary & 5{,}933{,}779 & 68.1 & 72.0 & -5.5 & -8.2 \\
Sweden & 7{,}121{,}754 & 81.8 & 85.0 & -3.7 & -15.1 \\
Spain & 17{,}937{,}712 & 167.3 & 173.6 & -3.6 & -8.3 \\
Germany & 67{,}146{,}981 & 765.8 & 793.7 & -3.5 & -9.6 \\
Croatia & 2{,}908{,}621 & 17.8 & 18.2 & -2.5 & -11.1 \\
Ireland & 3{,}580{,}975 & 35.1 & 34.8 & +0.6 & -7.0 \\
Slovakia & 3{,}993{,}065 & 40.7 & 39.8 & +2.2 & -6.5 \\
Cyprus & 890{,}693 & 3.3 & 3.1 & +3.3 & -4.8 \\
Portugal & 6{,}373{,}919 & 21.7 & 21.0 & +3.6 & +2.1 \\
Greece & 5{,}971{,}312 & 63.0 & 60.6 & +4.0 & +1.3 \\
Lithuania & 2{,}525{,}315 & 18.1 & 17.3 & +4.8 & -5.1 \\
Italy & 25{,}577{,}441 & 508.9 & 482.0 & +5.6 & +1.0 \\
Poland & 21{,}993{,}388 & 279.9 & 257.9 & +8.5 & -5.6 \\
Switzerland & 3{,}000{,}960 & 86.8 & 77.9 & +11.4 & +0.4 \\
France & 52{,}610{,}604 & 573.8 & 515.1 & +11.4 & +1.8 \\
Belgium & 9{,}170{,}389 & 132.8 & 111.8 & +18.7 & +7.5 \\
Latvia & 1{,}368{,}162 & 21.9 & 18.2 & +20.0 & +18.0 \\
Romania & 12{,}645{,}080 & 99.1 & 82.0 & +20.9 & +12.7 \\
Bulgaria & 4{,}087{,}643 & 35.6 & 29.0 & +22.7 & +16.9 \\
Slovenia & 1{,}430{,}391 & 17.4 & 14.0 & +24.5 & +20.1 \\
Malta & 143{,}277 & 1.0 & 0.7 & +35.6 & +9.6 \\
\midrule
\textbf{EU+UK+CH} & 317{,}590{,}992 & 3827.5 & 3858.8 & -0.8 & -8.3 \\
\bottomrule
\end{tabularx}
\smallskip
\raggedright\small
Bottom-up estimates from EUBUCCO v0.2 building footprints and TABULA per-vintage intensities; reported as-is, without calibration. Hotmaps figures are the all-class NUTS3 totals for 2015. A positive deviation means the bottom-up build is above the benchmark. Of the 29 countries, 19 fall within $\pm$15\,\% and 28 within $\pm$25\,\% on the raw deviation (only Malta, the smallest stock, exceeds the band), and the EU+UK+CH aggregate deviates by $-0.8$\,\%; the vintage-matched column shifts each country down by the 2015 to 2025 demand trend and is reported so the comparison is vintage-consistent. Three illustrative rows (Luxembourg, France, Finland) are discussed in the text.
\\[2pt]{\footnotesize\textit{Source: Authors' contribution.}}
\end{table}


\begin{figure}[H]
    \centering
    \caption{Validation of the bottom-up build against the Hotmaps baseline. Panel~(a)
    plots each country's bottom-up demand against its Hotmaps benchmark on log axes, with
    the dashed line parity and the shaded bands $\pm$15 and $\pm$25~per~cent. Panel~(b)
    gives the same deviation per country as a sorted bar chart, coloured green inside
    $\pm$15~per~cent, amber to $\pm$25 and red beyond.}
    \includegraphics[width=\textwidth]{paper/P03_validation.pdf}
    \label{fig:bottomup}
    \\[2pt]{\footnotesize\textit{Source: Authors' contribution.}}
\end{figure}


\subsection{Peak Dispatch, Peaking Investment, and Network Capital}
\label{sec:method:merit}

Levelised costs compare technologies as base-load suppliers of annual heat.
Hydrogen's most favourable case, however, is as a \emph{peaking} producer serving
only the cold-snap slice of the load-duration curve, the energy above the
2000-hour demand level, $\approx$12 to 18\% of annual heat
(Figure~\ref{fig:meritorder}a), when air-source heat-pump performance collapses and
the power system is stressed. We therefore model the winter peak explicitly as a
merit-order dispatch among the three quick-start options. Their marginal costs per
MWh of useful heat are
%
\begin{equation}
\label{eq:peakers}
    c^{\mathrm{gas}}_{c}
        = \frac{p^{\mathrm{gas}}_{c} + \varphi_{y}\, e_{\mathrm{g}}}{\eta_{\mathrm{g}}},
    \qquad
    c^{\mathrm{H_2}}_{c}
        = \frac{p^{\mathrm{H_2}}_{c}}{\eta_{\mathrm{h}}} + \sigma_{c},
    \qquad
    c^{\mathrm{HP}}_{c}
        = \frac{\pi^{\mathrm{peak}}_{c}}{\mathrm{COP}^{\mathrm{cold}}_{c}},
\end{equation}
%
where $\varphi_{y}$ is the carbon price, $e_{\mathrm{g}}$ the gas emission factor
(0.202~tCO$_2$/MWh), $\eta_{\mathrm{g}}{=}0.92$ and $\eta_{\mathrm{h}}{=}0.90$ the
boiler efficiencies, $\sigma_{c}$ the seasonal-storage adder ($\approx$\euro{}50/MWh-heat with domestic salt
caverns, $\approx$\euro{}100 without, at rates consistent in order with the bulk-storage
literature \citep{Lord2014GeologicStorage,Papadias2021BulkStorage,Talukdar2024UHSEurope} and applied on the
cavern geology mapped by \citet{Caglayan2020caverns}), and $\mathrm{COP}^{\mathrm{cold}}_{c} =
\max(0.60\,\mathrm{SCOP}_{c},\,1.8)$ the single country-level cold-snap heat-pump
coefficient of performance. The factor $0.60$ is the seasonal-to-cold-snap derate
calibrated to the low-temperature field measurements of \citet{Gibb2023Coming}. Across the cold-snap
design window an air-source unit delivers about three-fifths of its seasonal
coefficient, floored at $1.8$ for the coldest, highest-flow-temperature systems. The
seasonal coefficients themselves, $\mathrm{SCOP}_{c}$ of about 2.8 to 4.5 across the
stock, sit at the modern variable-speed end of the field evidence, consistent with
recent cold-climate and market-tracking data~\citep{Gibb2023Coming,EHPA2024market}; the
cold-snap derate above then keeps peak-hour dispatch from being credited with these
mild-weather values.

A third coefficient sizes the electricity network. The reinforcement adder of
Eq.~\eqref{eq:eladder} and the infrastructure bill of Eq.~\eqref{eq:infrabill} divide
by a flat $\mathrm{COP}^{\mathrm{el}}=2.3$ applied to every country, where
$\mathrm{COP}^{\mathrm{cold}}_{c}$ runs from 1.80 to 2.84 with a median of 2.04. The
flat value therefore sits above the median and understates the reinforcement both
carriers' comparison rests on, by about a tenth on the network bill, in the direction
that favours the heat pump. Correcting it is the first item for the next model run.
This single $\mathrm{COP}^{\mathrm{cold}}_{c}$ is used wherever the building or grid
peaker is priced, namely the merit-order dispatch of Eq.~\eqref{eq:peakers}, the
missing-money rent of Eq.~\eqref{eq:missingmoney} and the grid-reinforcement term of
Eq.~\eqref{eq:infrabill}; the deeper hour-resolved temperature derate that the 8760-hour
dispatch applies within the very coldest Dunkelflaute hours
(Appendix~\ref{sec:method:merit}, `Explicit hourly dispatch') is a distinct, sub-peak
quantity and is labelled as such where it is introduced, so that no two different cold
coefficients enter the same calculation. The scarcity electricity price
$\pi^{\mathrm{peak}}_{c}$ is derived from the power sector's own merit order. The
marginal winter unit, a gas OCGT versus a hydrogen-ready turbine at wholesale fuel
prices, sets it, plus a scarcity premium $\pi^{\mathrm{s}}$ spanning storage and
demand-response displacement,
%
\begin{equation}
\label{eq:peakprice}
    \pi^{\mathrm{peak}}_{c}
        = \min\!\left(
            \frac{p^{\mathrm{gas,w}} + \varphi_{y}\, e_{\mathrm{g}}}{\eta_{\mathrm{OCGT}}},\;
            \frac{p^{\mathrm{H_2,bb}}_{c}}{\eta_{\mathrm{T}}} + \sigma^{e}_{c}
          \right) + \pi^{\mathrm{s}},
\end{equation}
%
with turbine efficiencies $\eta_{\mathrm{OCGT}}{=}\eta_{\mathrm{T}}{=}0.40$, backbone
(not retail) hydrogen prices, and $\pi^{\mathrm{s}}\in\{30,60,120\}$~\euro{}/MWh
(low/central/high). The premium bounds the intra-seasonal flexibility cost that sits
above the marginal peaker's own short-run cost, that is, the value of the storage and
demand-response capacity that is displaced when the cold snap exhausts the cheaper
flexibility. The low value (\euro{}30/MWh) represents a deeply interconnected system
with abundant fast-responding demand response, the central value (\euro{}60/MWh) a
typical continental system, and the high value (\euro{}120/MWh) an islanded system with
limited battery storage in which scarcity rents approach the regulated value of lost
load. The three values bracket the plausible range without estimating a point. They are
stylised values we set, and no market has been observed to pay them, and we report them
as such. Two consequences follow from where the premium enters.

It cancels exactly from the hydrogen-versus-gas comparison in the power arena, where
both turbines buy at the same clearing price, so the power-arena win count holds at
seven across the whole \euro{}30 to 120/MWh-e band. A winner's rent is the premium times
its run hours, so capital recovery is exactly proportional to it. The 8.7~per~cent capacity-weighted recovery
computed at the central value is 4.4~per~cent at the low and 17.4 at the high, none of
which reaches the 100~per~cent unaided financeability requires.

The building peak behaves differently, because it is precisely this scarcity price that
the cold-snap heat pump is charged at, so there a higher $\pi^{\mathrm{s}}$ widens
hydrogen's niche and a lower one narrows it; this is the classic scarcity-pricing margin of low-utilisation capacity
\citep{Joskow2008capacity}. The sensitivity of hydrogen's merit-order position to
$\pi^{\mathrm{s}}$ is reported in full in the Results
(Section~\ref{sec:results:meritorder}, the winning-country count by scenario) and the
dedicated stress test of Appendix~\ref{sec:sensitivity:reinforce}; the missing-money
verdict (Eq.~\eqref{eq:missingmoney}) is invariant across the band, because it is the
peaker's running hours, and not the price level, that fail the investment test.
Hydrogen's heating share in each scenario is the exogenous adoption setting. The peak slice it would win under Eq.~\eqref{eq:peakers} is derived separately and compared with that setting, and never sets it.

Winning the dispatch is necessary but not sufficient. An investor must also recover the
peaker's capital from the hours it runs. The hydrogen peaker earns the inframarginal
rent against the cheapest alternative for the $H_{c}\approx500$ full-load hours of the
peak slice. This $500$-hour figure is not a free assumption but the duration of the
slice itself. It is the slice's energy divided by its capacity, a utilisation read from
the heat load-duration curve (Figure~\ref{fig:meritorder}a). It is not the count of
hours above the 2{,}000-hour threshold, which is 2{,}000 by construction. The slice
covers residential space heating only; industrial and process demand are outside the
model's scope and do not enter $H_{c}$.

This dwelling-level three-way contest between a gas boiler, an air-source heat pump and
a hydrogen boiler is distinct from the power-sector peaker stack of
Eq.~\eqref{eq:peakprice} and the hourly dispatch of Eq.~\eqref{eq:dispatch}. The two are
linked, because the heat-pump branch is priced at the endogenous scarcity electricity
price $\pi^{\mathrm{peak}}_{c}$ that the power-sector merit order sets, but the
electricity peaker that backs the grid runs far fewer full-load hours than the building
heat peaker, since the electricity peak is sharper and briefer than the heat peak. The
rent earns an annual margin per kW of capacity of
%
\begin{equation}
\label{eq:missingmoney}
    R_{c} = \left[\min\left(c^{\mathrm{gas}}_{c},\, c^{\mathrm{HP}}_{c}\right)
            - c^{\mathrm{H_2}}_{c}\right] \cdot \frac{H_{c}}{1000},
    \qquad
    \text{build iff } R_{c} \geq
        \bigl[\mathrm{CRF}(r_{c}, n) + m\bigr]\, K,
\end{equation}
%
where $K$ is the hydrogen-boiler CAPEX (\euro{}600/kW), $n{=}20$ years, $m{=}2\%$ fixed
O\&M, and $r_{c}$ the country WACC.

The power-sector peaker is a different asset and takes its own condition, which every
recovery figure in Section~\ref{sec:power:caveat} and Appendix~\ref{sec:results:dispatch}
uses. Where hydrogen is the cheaper firm unit it sets the clearing price at its own
short-run cost plus the scarcity premium $\pi^{\mathrm{s}}$, so its rent is that premium
over the hours it runs and its capital recovery is
%
\begin{equation}
\label{eq:powerrecovery}
    R^{\mathrm{pw}}_{c}
      = \frac{\pi^{\mathrm{s}}\,\mathrm{FLH}_{c}/1000}
             {\bigl[\mathrm{CRF}(r_{c},\,25) + f\bigr]\, K^{\mathrm{T}}},
    \qquad f = 0.025,
\end{equation}
%

with $\mathrm{FLH}_{c}$ the peaker's full-load hours, $f$ a fixed operating charge of
2.5~per~cent a year, and $K^{\mathrm{T}}$ the hydrogen-ready turbine's overnight capital
cost of \euro{}600/kW, against \euro{}450/kW for a conventional open-cycle machine. The
turbine's 25-year life and its 2.5~per~cent fixed operating charge differ from the
boiler's 20 years and 2~per~cent above, and the two assets share the \euro{}600/kW figure
by coincidence rather than by construction, so the two conditions are not interchangeable.
The denominator is the annualised charge quoted as \euro{}51 to 66/kW-yr for the five
markets that build: Denmark's $36.40+15.00=51.40$ and Poland's $51.49+15.00=66.49$.
Section~\ref{sec:results:missingmoney} shows that the heat-side build condition fails
everywhere, and the magnitude of the failure is large. Over the roughly $500$ full-load hours of the slice the inframarginal rent recovers only
a minority of the heating-side peaking boiler's annualised capital, with the best cases (Sweden and Denmark,
under H2 Push) at approximately 40~per~cent, and 12 of the 29 winning country-scenarios
recovering under 20~per~cent, because the inframarginal rent a peaker earns over so few
hours falls far short of an annualised capital requirement that runs from \euro{}52 to
73/kW-yr across the 29 markets, and from \euro{}52 to 71 across the winning ones.

Resolved to the hour and accumulated across the milestone years to 2050, the same
shortfall is about \euro{}357~bn for the hydrogen peaker fleet EU-wide under the H2~Push
scenario, the central estimate carried in Section~\ref{sec:power:caveat}, and sweeps to
roughly \euro{}379 to 563~bn across the remaining scenarios (\euro{}379~bn Net Zero,
\euro{}418~bn Stated Policies, \euro{}563~bn Current Policies; against about \euro{}130
to 142~bn for the gas fleet), so the gap is structural and persists under an explicit
8760-hour treatment; the full per-country-scenario detail is in
Section~\ref{sec:results:missingmoney} and Appendix~\ref{sec:results:dispatch}.

Finally, each pathway's network capital requirement aggregates three components
from the model's own technology shifts,
%
\begin{equation}
\label{eq:infrabill}
    B_{c} = \underbrace{\Delta S^{\mathrm{HP}}_{c}\,
            \frac{P^{\mathrm{pk}}_{c}}{\mathrm{COP}^{\mathrm{el}}}\,
            \kappa\, u^{\mathrm{el}}}_{\text{grid reinforcement}}
          + \underbrace{\Delta S^{\mathrm{DH}}_{c}\, N_{c}\,
            u^{\mathrm{DH}}}_{\text{DH expansion}}
          + \underbrace{S^{\mathrm{H_2}}_{c}\, N_{c}
            \bigl[\gamma_{c} u^{\mathrm{conv}} + (1-\gamma_{c})
            u^{\mathrm{new}}\bigr]}_{\text{H}_2\text{ network}},
\end{equation}
%
where $\Delta S$ are the 2025$\to$2050 share shifts, $P^{\mathrm{pk}}_{c}$ the heat
peak, $\kappa{=}0.5$ the coincidence factor,
$u^{\mathrm{el}}\in\{400,800,1600\}$~\euro{}/kW the reinforcement cost band, $N_{c}$
the dwelling stock, $u^{\mathrm{DH}}\in\{3000,5000,8000\}$~\euro{} per connection,
$\gamma_{c}$ the gas-grid coverage, and $u^{\mathrm{conv}}/u^{\mathrm{new}}$ the
hydrogen conversion/new-build connection cost bands of
Appendix~\ref{sec:results:h2infra}.

\paragraph{Explicit hourly dispatch.} Equations~\eqref{eq:peakers} to
\eqref{eq:missingmoney} treat the peak as a single load-duration slice. To make the
cold-snap hours and the peaker's utilisation explicit, and to test whether the
missing-money result survives an hour-resolved treatment, we add an 8760-hour dispatch
of the power sector against a per-country installed-capacity stack built from national
2050 strategies and NECPs and cross-checked against the ENTSO-E TYNDP~2024 (data note in
the supplementary material). Hourly electricity demand $D_{c,t}$ combines a calibrated
national baseline-load shape with heat-pump electricity
$Q^{\mathrm{heat}}_{c,t}/\mathrm{COP}_{c,t}$, where the hourly COP $\mathrm{COP}_{c,t}$
is \emph{temperature-dependent}, tracking the seasonal coefficient in mild hours and
falling towards a sub-peak floor of $0.55\,\overline{\mathrm{SCOP}}$ in the very coldest hours
of the synthetic Dunkelflaute, so the electricity peak is sharper than the heat peak,
the physical reason a cold, still, dark hour stresses the grid.

The overbar marks a study-area mean. The hourly layer blends the air-source and
ground-source seasonal coefficients three to one and applies the result to every
country, without the climate multiplier (0.85 in Finland to 1.35 in Cyprus) that the
levelised-cost layer applies through $\mathrm{SCOP}_{c}$. Capacity-weighted across the
peaker fleet that multiplier averages 0.966, so the per-country split of the firm fleet
carries a bias of a few per cent and the coldest markets are the ones understated. This
hour-resolved floor ($0.55\,\overline{\mathrm{SCOP}}$) is deliberately deeper than, and
distinct from, the single country-level cold-snap coefficient
$\mathrm{COP}^{\mathrm{cold}}_{c}=\max(0.60\,\mathrm{SCOP}_{c},1.8)$ of
Eq.~\eqref{eq:peakers}. The latter is the building or grid peaker's design value over
the whole cold-snap window and prices the merit order and the missing-money rent,
whereas the $0.55$ floor applies only to the handful of extreme hours that set the
electricity peak in this dispatch layer. The two are used in different calculations and
are never substituted for one another, so the merit-order position of
Eq.~\eqref{eq:peakers} and the rent of Eq.~\eqref{eq:missingmoney} are computed at
$\mathrm{COP}^{\mathrm{cold}}_{c}$ throughout. Table~\ref{tab:copmap} states the mapping. The three derates are distinct coefficients entering distinct equations, and none is used in another's place.

\begin{table}[H]
    \centering
    \caption{The three cold-weather heat-pump derates and where each enters.}
    \label{tab:copmap}
    \small
    \begin{tabular}{@{}p{0.24\textwidth}p{0.40\textwidth}p{0.30\textwidth}@{}}
        \toprule
        Derate & Definition & Enters \\
        \midrule
        Merit-order / missing-money
            & $\mathrm{COP}^{\mathrm{cold}}_{c}=\max(0.60\,\mathrm{SCOP}_{c},\,1.8)$ & Eqs.~\eqref{eq:peakers}, \eqref{eq:missingmoney} \\ Network reinforcement & $\mathrm{COP}^{\mathrm{el}}=2.3$, flat across countries & Eqs.~\eqref{eq:eladder}, \eqref{eq:infrabill} \\
        Hourly extreme-hour dispatch
            & $\mathrm{COP}_{t}\to 0.55\,\overline{\mathrm{SCOP}}$ (coldest hours)
            & Demand $D_{c,t}$ in Eq.~\eqref{eq:dispatch} \\
        \bottomrule
    \end{tabular}
    \\[2pt]{\footnotesize\textit{Source: Authors' contribution.}}
\end{table}

Each generator $g$ supplies
$a_{g,c,t}\,P_{g,c}$, where $a$ is a calibrated hourly availability (solar diurnal
to seasonal; wind an AR(1) process with a forced cold-snap \emph{Dunkelflaute};
hydro run-of-river seasonal; nuclear flat must-run), and a flexibility budget
$f_{c,t}$ (a four-hour battery plus the hydro reservoir) shaves the daily peak. The
residual served by the dispatchable fleet, and the peaker energy and capacity above
the firm gas fleet $P^{\mathrm{gas}}_{c}$, are
%
\begin{equation}
\label{eq:dispatch}
    r_{c,t} = D_{c,t} - \!\!\sum_{g\in\{\mathrm{VRE,nuc,bio}\}}\!\!\! a_{g,c,t}\,P_{g,c} - f_{c,t},
    \quad
    E^{\mathrm{pk}}_{c} = \sum_{t}\max\!\bigl(0,\, r_{c,t}-P^{\mathrm{gas}}_{c}\bigr),
    \quad
    G^{\mathrm{pk}}_{c} = \max_{t}\,r_{c,t}.
\end{equation}
%
The peaker sized to $G^{\mathrm{pk}}_{c}$ runs $E^{\mathrm{pk}}_{c}/G^{\mathrm{pk}}_{c}$
full-load hours; we price it as hydrogen- versus gas-fuelled with full economics, and
all inputs evolve over time, among them the \emph{rising} carbon price on gas, a
\emph{declining} gas wholesale price (\euro{}38$\to$28/MWh) and hydrogen price, an
\emph{improving} turbine efficiency (40\% in 2025, rising to 42\% for gas and 44\% for
the hydrogen turbine by 2050, or 48\% under H2~Push, which assumes an advanced
simple-cycle machine, the latter slightly faster, reflecting H$_2$-combustor R\&D), a
declining seasonal-storage cost, variable O\&M, country labour, and annualised CAPEX
(\euro{}600/kW H$_2$-ready turbine, \euro{}450/kW OCGT, 25-year life).

These turbine efficiencies sit within the open-cycle range of the Danish Energy Agency
technology catalogue, in which advanced and hydrogen-ready simple-cycle units reach the
mid-to-high forties per cent by 2050~\citep{DEA2023technology}, so the peaking economics
do not rest on optimistic turbine performance. We accumulate profit and system cost
across the 2025 to 2050 milestone years. The structural link to the buildings model is
the endogenous peak price $\pi^{\mathrm{peak}}_{c}$ of Eq.~\eqref{eq:peakprice}. We
bound the three assumptions to which the peaker requirement is most sensitive:
\emph{interconnection} (the islanded sum of national peakers versus an EU copper-plate
fleet sized to the aggregate residual peak), the \emph{depth} of the synthetic cold-snap
\emph{Dunkelflaute} (cold-snap wind swept from 15\% to 40\% of normal), and the
\emph{length} of that cold-and-still event (swept from two days to ten). The resulting
ranges are reported with the result.

The third is the largest of the three. The event length is a
hard-coded six days, and both the heat spike and the wind lull run on the same window in
every country at once, so the other two sweeps leave it fixed. Sweeping it moves the
fleet sized on the inflexible load from 253.7~GW at two days to 302.0~GW at ten, against 277.7 at the
central six, which is a wider band than either of the other two produces. It also
disposes of a natural explanation for the modest reduction a partially
flexible heat-pump load delivers. That reduction might be attributed to the event being
multi-day, on the
argument that intra-day pre-heating cannot move load out of a six-day window. That argument restates the assumption, and the sweep refutes it; the
reduction is nearly invariant across the range, at 8.9 to 11.6~per~cent, because what binds is the share of load that can shift, and the length of the event does not. A fourth
assumption, that the heat-pump load is completely inflexible, is bounded separately at
4.3 to 10.7~per~cent of the fleet across three energy-conserving smoothing operators.

\paragraph{Unit-commitment dispatch of the firm fleet.} The reduced-form merit order of
Eq.~\eqref{eq:dispatch} sizes the dispatchable layer by a simple threshold on the
residual, which treats the firm fleet as perfectly flexible. To remove that
idealisation, and to address the most material structural concern about the power
coupling, the dispatchable layer is now resolved by an explicit \emph{unit-commitment}
optimisation.

The dispatchable fleet is split into a firm combined-cycle (CCGT) class and fast
open-cycle gas or hydrogen-turbine peakers, and for each country and milestone year the
model minimises the weighted fuel, carbon and variable operating cost plus start-up
cost, \[ \min
\sum_{d}\omega_d\sum_{h}\left[\sum_{g}\left(c^{\mathrm{var}}_{g}\,p_{g,d,h} +
c^{\mathrm{su}}_{g}\,v_{g,d,h}\right) + \pi^{\mathrm{voll}}\,u^{\mathrm{ns}}_{d,h} +
\pi^{\mathrm{res}}\,s^{\mathrm{res}}_{d,h}\right], \]

subject to commitment and physical
constraints on each class: minimum stable output ($p^{\min}_{g}\,n_{g,d,h}\le
p_{g,d,h}\le \bar p_{g}\,n_{g,d,h}$), ramp limits, minimum up and down times, a 5~per
cent operating-reserve margin on the fixed hourly net load (committed capacity must
exceed served load plus the margin), and the battery and reservoir state of charge
carried hour to hour, with $\omega_d$ the weight of representative day $d$, $n_{g,d,h}$
the count of committed $\approx$1~GW units, $v_{g,d,h}$ start-ups, and
$u^{\mathrm{ns}},s^{\mathrm{res}}$ slacks for unserved energy (value of lost load
$\pi^{\mathrm{voll}}=$~\euro{}8{,}000/MWh) and reserve shortfall
($\pi^{\mathrm{res}}=$~\euro{}1{,}000/MW-h). The optimisation reuses the reduced form's
demand, renewable availability (including the synthetic \emph{Dunkelflaute}),
temperature-dependent heat-pump COP derate and causal storage dispatch verbatim, so it
differs from the reduced-form merit order \emph{only} in how the dispatchable fleet
is committed.

The upgrade therefore removes the fleet-flexibility idealisation,
that the firm fleet could follow the residual perfectly, but it retains the
reduced form's calibrated variable-renewable availability profiles and its causal,
non-co-optimised storage heuristic; both remain assumptions of this layer and are
carried as such in the limitations (Appendix~\ref{sec:limitations}).

The 8760 hours are reduced to a weighted set of representative
days per country, force-including as weight-one singletons the top peaker-signal days
that together hold at least 99~per cent of annual peaker energy (so the convex
peaker quantity is not averaged away), and clustering the remaining low-peaker days;
the per-country problems are solved in parallel. Techno-economics follow the Danish
Energy Agency Technology Catalogue 2023 \citep{DEA2023technology} and the ENTSO-E
TYNDP~2024: a CCGT minimum stable output of 0.30 of capacity, a CCGT ramp of 2.4 per
hour, four-hour minimum up and three-hour minimum down times, against faster
open-cycle and hydrogen peakers (minimum stable output 0.20, ramp 4.0). The CCGT
ramp and minimum-stable-output values, and the 5~per cent operating-reserve
standard, are judgement-call inputs carried as assumptions and tested in the
limitations (Appendix~\ref{sec:limitations}).

The production model is a linear
relaxation of the commitment problem; against the mixed-integer formulation on spot
countries the integrality gap is only 0.04 to 0.11~per cent, so commitment
integrality is immaterial to the result and the relaxation is used throughout. The
former reduced-form merit order of Eq.~\eqref{eq:dispatch} is retained as a validation
cross-check, run as a limit test with the commitment constraints disabled and
reported with the rest of the hourly-dispatch results in
Appendix~\ref{sec:results:dispatch}. The full specification, the representative-day
reducer and the validation gates are documented in the unit-commitment methods note
accompanying the supplementary material.


\subsection{Uncertainty Quantification}
\label{sec:method:uncertainty}

Parametric uncertainty is characterised with a Monte Carlo approach using $N = 200$
samples per scenario. The following parameters are sampled stochastically: the
per-country envelope-design intensity decline rate
$r_{c,s}\sim\mathcal{N}(\mu_{c,s},\sigma_{c,s})$, with
$\sigma_{c,s}\approx0.21\mu_{c,s}$ so that the published low/high bounds span the
5th to 95th percentile; the technology share targets for 2050; the heat-pump
ground/air split; and the fuel-price multipliers (log-normal with $\sigma = 0.15$).
Because renovation pace is driven partly by common EU policy (the EPBD and the
Renovation Wave) and partly by national factors, the per-country rate shocks are
correlated through a shared factor:
$z_c=\sqrt{\rho}\,z_{\text{EU}}+\sqrt{1-\rho}\,z_c^{\text{idio}}$ with $\rho=0.5$.
This preserves each country's marginal distribution (so per-country bands and all
medians are unchanged) while preventing the EU-aggregate band from being
artificially narrowed by cross-country diversification; the resulting EU band is
about 2.2 times wider than under independent draws. The parameter $\rho$ is a
structural knob (Appendix~\ref{sec:limitations}). Results are reported as the 10th,
50th and 90th percentile quantiles across samples.

The choice of $N = 200$ samples is supported by a convergence test, and we report what
it does and does not establish. The running median of the headline metric, 2050 EU
buildings CO$_2$ emissions, settles inside $\pm 1\,\%$ of its final value after 41 draws under Current Policies and 48 under Stated Policies, after 96 under Net~Zero, whose two-branch mixture converges slowly, and only after 139 under H2~Push, the slowest of the four. The deciles are slower
still and are not converged at $N = 200$. The running 10th percentile last leaves that
tube between draws 156 and 188 in all four scenarios (Figure~\ref{fig:mcconv}). A running
quantile can re-enter the tube and leave it again, so each figure above is the last draw
at which the series is still outside, which is the conservative of the two readings.
Medians are therefore reliable at this sample size and decile bands should be read as
indicative.

\begin{figure}[H]
    \centering
    \caption{Running-quantile convergence of the 2050 buildings CO$_2$ metric. Each
    series is the quantile recomputed on the first $k$ draws, expressed as a deviation
    from its value at $N = 200$. The shaded band is the $\pm 1\,\%$ tube and the marker
    on each series is the last draw at which it sits outside. Panel~(b)'s first few dozen
    draws overshoot to $+27\,\%$ and are clipped, so the axis resolves the settled tail
    the claim is about.}
    \includegraphics[width=\textwidth]{paper/P42_mc_convergence.pdf}
    \label{fig:mcconv}
    \\[2pt]{\footnotesize\textit{Source: Authors' contribution.}}
\end{figure}


\subsection{The Anatomy of the Cost Ordering}
\label{sec:results:lcohanatomy}
We decompose these levelised costs into their components, to show where the
ordering comes from. Figure~\ref{fig:fig7} resolves the gas boiler, the air-source
heat pump and the hydrogen boiler into capital, fixed and variable operation and
maintenance, fuel, ETS2 carbon, and delivery and storage. On the EU median stacks,
which total \euro{}129.3/MWh for the gas boiler, \euro{}111.5/MWh for the
air-source heat pump and \euro{}124.8/MWh for the hydrogen boiler, the first
reading is that fuel is the largest single component of every one of them, at
45.3~per cent of the gas-boiler stack, 54.2~per cent of the heat-pump stack and
43.6~per cent of the hydrogen-boiler stack, so it is the fuel- and
electricity-price trajectories of Figure~\ref{fig:fig5}a that govern long-run
competitiveness in all three cases, the heat pump included.

Three structural features then separate them. The first is carbon. The ETS2 wedge sits
on the gas stack alone, at 33.2~per cent of its cost, and it is what holds gas at the
top of the 2050 ordering. The second is capital, which runs the opposite way from the
usual intuition. The heat pump is the \emph{most} capital-heavy of the three, at
27.7~per cent of its levelised cost against 17.2~per cent for the hydrogen boiler and
11.8~per cent for the gas boiler, which is why the household-level barrier to adoption
is an up-front financing one even though fuel remains its single largest cost component.
The third, and the decisive one for the comparison this paper turns on, is what it costs
to move the molecule and to hold it.

Delivery and seasonal storage together take 27.0~per cent of the hydrogen stack, just
under the 27.8~per cent its capital and fixed operation and maintenance take together,
against 4.4~per cent of the heat-pump stack for network reinforcement and nothing at all
for gas, whose distribution grid is already built and already inside the retail price.
This is the hydrogen boiler's largest structural burden. The three stacks therefore
describe an efficient machine that is expensive to buy and runs on expensive taxed
electricity, set against a cheap molecule that is expensive to deliver and expensive to
keep from the summer in which it is made to the winter in which it is burned, and the
ordering between them turns on the balance of those burdens, with no single dominant
item.

\begin{figure}[H]
    \centering
    \caption{The anatomy of the levelised cost of heat by technology, 2030 and 2050 under
    Stated~Policies. Each component is a median taken independently across the 29 markets
    and the bars are their sum, so the stacked totals of about \euro{}112, \euro{}125 and
    \euro{}129 differ from the \euro{}118.5, \euro{}122.8 and \euro{}132.7 per-country
    median costs quoted throughout the text. The figure is for the composition, not the
    level.}
    \includegraphics[width=0.86\textwidth]{paper/P07_lcoh_anatomy.pdf}
    \label{fig:fig7}
    \\[2pt]{\footnotesize\textit{Source: Authors' contribution.}}
\end{figure}

%── A.5 CO2 Emissions ─────────────────────────────────────────────────────────

\subsection{Decomposition: Demand, Efficiency and Grid Cleanness}
\label{sec:results:decomp}

The headline emissions fall is driven by what happens
on the supply side. Table~\ref{tab:threelayer} separates the Stated Policies result
into three layers. Useful heat demand, what the building physically requires, falls
only 9 per cent, since envelope renovation barely outpaces stock growth.
\emph{Delivered} (final) energy falls 31 per cent, because the shift to heat pumps
(with a seasonal coefficient of performance of approximately 3) decouples delivered
energy from useful heat. The implied EU system efficiency rises from 0.97 in 2025
to 1.29 by 2050. Emissions then fall 64 per cent, as the carrier mix moves off gas
and oil and the grid behind the electrified share decarbonises. Decarbonisation is
therefore a supply-side phenomenon, driven by source switching, conversion efficiency and grid cleanness, with no contribution from demand destruction.

\begin{table}[H]
    \centering
    \caption{The three layers of the result (Stated Policies, 2025 to 2050).}
    \label{tab:threelayer}
    \begin{tabular}{lrrr}
        \toprule
        Layer & 2025 & 2050 & $\Delta$ \\
        \midrule
        Useful heat demand (TWh)  & 3{,}827 & 3{,}483 & $-9\%$ \\
        Final/delivered energy (TWh) & 3{,}928 & 2{,}695 & $-31\%$ \\
        CO$_2$ emissions (MtCO$_2$)  &   644 &   231 & $-64\%$ \\
        \bottomrule
    \end{tabular}
    \\[2pt]{\footnotesize\textit{Source: Authors' contribution.}}
\end{table}

The depth of the cut nonetheless depends in part on the exogenous grid trajectory,
and we test how much. This attribution runs on a single central draw in place of the 200-draw median. It shares the 644~MtCO$_2$ 2025 level with the ensemble, and reaches a Stated Policies 2050 level of about 264~MtCO$_2$ against the 231~MtCO$_2$ median reported elsewhere. The two are the same pathway read at different points of its own distribution, and only the split between the two causes is at issue here, so the percentages below are computed on the 644 to 264 fall of this draw and not on the median pair.

Recomputing that pathway with both grid and district-heat
carbon intensity frozen at their 2025 levels, while keeping the same technology
mix, leaves 2050 emissions under Stated Policies at approximately 407~MtCO$_2$,
still a 36 per cent reduction on 2025. Technology switching alone therefore
delivers roughly \emph{three-fifths} of the scenario's emissions fall (62 per cent),
achievable even on a 2025-vintage grid, with grid and district-heat decarbonisation
contributing the remaining two-fifths (38 per cent). The result is primarily a technology-switching outcome, and the assumed grid trajectory
is not what produces it; but the grid
assumption governs the remaining two-fifths, and the deep (98 per cent) outcome of Net Zero
is unattainable without it (Appendix~\ref{sec:limitations}).

%── A.6 Conditions for Hydrogen Deployment (RQ1) ──────────────────────────────

\subsection{Hydrogen Distribution-Infrastructure Cost}
\label{sec:results:h2infra}

The headline gap of Figure~\ref{fig:h2gap} charges each carrier for its own
last-mile network. Hydrogen pays a distribution adder set at the coverage-weighted
blend of conversion and new-build cost, at the central cost bound, and electric
heating pays the low- and medium-voltage reinforcement that its added coincident
winter peak requires (Appendix~\ref{sec:method:h2}). Free reuse of the existing gas
grid at no incremental cost is retained only as the generous end of the sweep
below, because a methane network cannot carry 100 per cent hydrogen unmodified.
Polyethylene mains are largely reusable, but cast-iron and ductile-iron mains,
meters, regulators and purging and leak-detection equipment require replacement or
upgrade, and ungridded areas need a network built from scratch
\citep{NorthernGasNetworks2018h21,GasForClimate2020backbone}. We therefore sweep
hydrogen's delivery cost across four per-country treatments, each annualised over a
40-year network life at the country WACC and spread over the useful heat a dwelling
consumes per year (climate-scaled around an EU average of approximately 13~MWh):
\begin{enumerate}
    \item \textbf{Free gas-grid reuse (lower bound).} Hydrogen inherits the
    existing distribution grid at no incremental cost. This is the free-reuse assumption.
    \item \textbf{Retrofit conversion.} The whole existing gas grid is
    \emph{converted} to 100 per cent hydrogen. This is still optimistic, since it
    assumes a usable grid already reaches every home and needs only conversion.
    \item \textbf{Coverage-weighted blend (the model default).} Conversion cost is
    charged where the gas grid already reaches the dwelling and new-build cost
    where it does not, weighted by each country's gas-grid coverage $\gamma_c$, as
    in the hydrogen term of Eq.~\eqref{eq:infrabill}.
    \item \textbf{New-build (upper bound).} A dedicated hydrogen distribution
    network is built from scratch to every dwelling.
\end{enumerate}
The per-connection capital cost is itself genuinely uncertain, so we carry each
treatment as a sensitivity band over a low, central and
high range from the conversion literature: conversion at
\euro{}1{,}000/1{,}500/3{,}000 per dwelling, and new-build at
\euro{}3{,}000/5{,}000/8{,}000 per dwelling.

Two things follow from the sweep, one about direction and one about geography. The
direction is monotone. Free gas-grid reuse is the most favourable
treatment for hydrogen and full new-build the least, with retrofit conversion and
the coverage-weighted blend in between, and the EU-mean gap between the hydrogen
boiler and the best available heat pump widens against hydrogen at every step along
that sequence. The blend is the treatment the headline of
Section~\ref{sec:results:h2conditions} reports, so the headline sits at the third of the four rungs.

The sweep has been recomputed on the storage-inclusive basis this paper uses, so
every arm now carries the seasonal store and all of them are comparable both with one
another and with the headline. The EU-mean hydrogen-minus-heat-pump gap runs from
\euro{}0.0/MWh under free gas-grid reuse, through $+$\euro{}8.0 under retrofit
conversion and $+$\euro{}19.7 under the coverage-weighted blend used as the default,
to $+$\euro{}26.9 under full new build, reaching $+$\euro{}43.1 at the new-build high
bound. The count of markets at or below parity for hydrogen falls from 12 to 9, 7, 3
and finally 2 across those treatments. Only Denmark and Germany survive the most
expensive one, Denmark running from $-$\euro{}44.3/MWh under free reuse to
$-$\euro{}28.0 under new build and $-$\euro{}18.2 at the high bound, and Germany from
$-$\euro{}33.2 to $-$\euro{}16.2 and $-$\euro{}6.0.

The earlier reading drawn from the pre-storage version of this sweep, that hydrogen's
base-load advantage survives every network treatment in a shrinking but never empty
set of countries, does not survive the recomputation. The set does shrink to almost
nothing, and outside Denmark and Germany hydrogen's base-load case depends on
inheriting distribution assets it does not pay to build. For scale, the seasonal-storage
term the earlier sweep omitted is \euro{}12.1/MWh of useful heat in a cavern market such
as Germany and \euro{}28.4/MWh in a cavern-free one such as Spain
(Appendix~\ref{sec:method:h2}), which is of the same order as the \euro{}27/MWh the four
network rungs span between them.

What the sweep does establish, and what does not depend on the level of the
base-load cost, is that network cost redraws the map along storage geology and never across it. The countries least affected by charging
hydrogen for its pipes are the cold, high-electricity-price north-west, where the
cold season is long enough to spread a network charge thinly over a large annual
heat demand. The warm south and the low-price east are affected most, because a
small annual heat demand carries the same per-connection capital over far fewer
megawatt-hours. Those are the same two groups the seasonal-storage charge separates,
so the two effects reinforce one another. Two qualifications belong with this
result. The gas-grid asset-utilisation co-benefit is real but only partial, since
even the legacy network is not free to repurpose. And what advantage hydrogen
retains is a fiscal artefact as much as a technology one, because the comparison
sets an untaxed molecule priced delivered against electricity carrying
the full retail levy stack; that asymmetry, and the sensitivity that would test it,
is set out in Appendix~\ref{sec:limitations}.

%── A.7 Merit order: hydrogen as a peaking producer ───────────────────────────

\subsection{Hourly Dispatch, Capacity-Short Hours and the Peaker's Capital}
\label{sec:results:dispatch}

This completes Test~3 by resolving its single load-duration slice to the hour. The
missing-money result above rests on that one slice; resolved to the hour it holds and
sharpens. Figure~\ref{fig:dispatch} makes the point that peaker utilisation is small. It
is the one-day hourly snapshot of the capacity-short hours. As a robustness check we
dispatch an 8,760-hour power system against a per-country capacity stack
(Eq.~\eqref{eq:dispatch}). The model reproduces the standard pattern
(Figure~\ref{fig:dispatch}). On a cold, still, dark winter day the residual demand after
must-run and renewable generation spikes in the morning and again in the evening,
exactly when electrified-heat demand peaks, solar is absent and wind has collapsed into
a \emph{Dunkelflaute}. These are the hours in which the system runs short and a firm
peaker must fill the gap. Two features matter for hydrogen.

First, the scarcity is a \emph{system-adequacy} phenomenon, and heating alone does not
create it. Heat-pump electricity \emph{alone} never exhausts the generation fleet in any
country-year we model; it is the \emph{sum} of baseline load and electrified heat, set
against a low-VRE cold snap, that creates the gap. Second, the peaker's utilisation is
small and falls as renewables and storage build out, from over a thousand full-load
hours in the gas-dependent 2025 system to about 139 by 2050, spread across a mean of
roughly 440 capacity-short hours in the markets that build a peaker, squarely in the pure-peaking regime that the TYNDP~2024
scenarios assign to the EU gas and hydrogen-turbine fleet. The same scarcity is sharply
seasonal. Aggregated to daily means across the year (Figure~\ref{fig:seasonality}), the
firm-capacity deficit concentrates in the winter months and all but vanishes in summer,
so the peaker earns its keep in a narrow cold-season window.

\begin{figure}[H]
    \centering
    \caption{Cold-snap dispatch and the peaker gap (EU-27+UK+CH, 2050).}
    \includegraphics[width=\textwidth]{paper/P20_dispatch_coldsnap.pdf}
    \label{fig:dispatch}
    \\[2pt]{\footnotesize\textit{Source: Authors' contribution.}}
\end{figure}

\begin{figure}[H]
    \centering
    \caption{Seasonality of the peaking need (EU-27+UK+CH, 2050, daily mean). The shaded
    bands are the two winter halves of the calendar year, December to February. In the
    upper panel the stacked supply runs above the demand line through the summer, which is
    surplus available capacity and not generation; the model does not curtail it explicitly.
    In the lower panel the total-peaker line is the sum of the two stacked areas beneath it.}
    \includegraphics[width=\textwidth]{paper/P25_dispatch_seasonality.pdf}
    \label{fig:seasonality}
    \\[2pt]{\footnotesize\textit{Source: Authors' contribution.}}
\end{figure}

Pricing that peaker with full economics confirms and sharpens the base-case result.
On operating cost, hydrogen undercuts the gas turbine only where the carbon price is
high or hydrogen is deliberately cheap. In 2050 the hydrogen peaker's variable cost
beats gas in 0 of 29 countries under Current and Stated Policies, but, on the
even-handed series holding both turbines at 0.40, in 7 of 29
under both Net Zero and H2~Push. In each the \euro{}350/tCO$_2$ carbon price lifts the
gas turbine to about \euro{}256/MWh, and under H2~Push the rapid hydrogen decline
lowers the hydrogen turbine beneath it in the same seven markets. The same carbon-driven pattern documented in the main text recurs
at the building and district-heating levels.

On the investment question, however, \textbf{both fuels miss the money,
cumulatively, to 2050} (Figure~\ref{fig:peakercum}). The cumulative loss is the
milestone-year accumulation of each fleet's annual profit shortfall defined with
Eq.~\eqref{eq:dispatch}, that is, the sum over the modelled years of the gas and
hydrogen peaker losses. Summed across the EU-27+UK+CH and
integrated over the milestone years, the gas peaker fleet loses about \euro{}142~bn and
the hydrogen peaker fleet about \euro{}357~bn under the central H2~Push scenario,
sweeping to \euro{}379 to 563~bn across the remaining scenarios. Those figures come from the reduced-form dispatch, which also sets the fleet and the run hours; the explicit
unit-commitment layer below runs the same economics on its own dispatch and gives
\euro{}358.4~bn with its constraints off and \euro{}356.0~bn with them on. The losses arise because the firm capacity that electrified heat requires for adequacy, about 278~GW EU-wide in the central islanded case and lower under interconnection (see below), runs far too few hours to repay itself from market revenue. The gas peaker recovers between about 7 and 20~per cent of its capital, the classic missing money of low-utilisation plant.

We size and dispatch this fleet against demand with the storage fleet's state of charge
tracked hour by hour, so that a multi-day \emph{Dunkelflaute} depletes the battery and
the firm peaker must cover the sustained deficit; batteries serve the daily peak and
cannot cover a week-long cold snap. The hydrogen peaker loses two to four times as much
as gas, driven by its higher fuel-plus-storage cost, and the worst case is Current
Policies, where stranded-expensive hydrogen never approaches viability. The system
itself stays adequate; modelled loss of load averages about two hours per year and never
exceeds the three-hour standard, so this is an economic verdict. The peaker requirement,
and hence the absolute totals, turns on two assumptions that we test explicitly.

\emph{Cross-border interconnection} lowers the requirement from the islanded case (no
trade, with each country self-supplying its peak) towards an EU copper-plate fleet sized
to the aggregate residual peak (perfect pooling), and pooling helps least in a deep,
continent-wide \emph{Dunkelflaute}, when every country is short at once, and most in
milder, regional cold snaps, so trade cannot be relied upon to remove the firm-capacity
need. The \emph{depth} of the modelled cold snap moves the requirement in both
directions around the central figure. Across cold-snap wind availabilities of 0.15, 0.25
and 0.40 the islanded requirement runs 291, 278 and 251~GW against a copper-plate 267,
239 and 205~GW, so the whole two-way grid spans 205 to 291~GW and the central islanded
case reproduces the 278~GW headline exactly. The euro losses scale with this capacity,
but the sign and ranking are invariant. In every case the hydrogen peaker loses two to
four times what the gas peaker does, and neither recovers its capital from market
revenue.

\paragraph{The verdict under explicit unit commitment.}
Resolving the dispatchable layer with the unit-commitment optimisation of
Appendix~\ref{sec:method:merit} confirms this result on a more realistic
representation of how the firm fleet runs. The validation is exact first. With the
commitment constraints disabled the optimisation reproduces the reduced-form merit
order to within 0.33~per cent per country on dispatched energy and 0.40~per cent on
peak capacity, and puts the fleet at 280.5~GW. The 2.8~GW difference is a Romanian peaker that the sizing pass builds and the reduced form does not. The unchanged peaker economics then feed through to a cumulative-2050 \emph{limit case}
of $-$\euro{}358.4~bn for hydrogen and $-$\euro{}143.2~bn for gas under H2~Push, within
\euro{}1.1 and \euro{}0.8~bn of the reduced-form headline of
Section~\ref{sec:power:caveat}, the residual being the Romanian peaker. The
unit-commitment layer is therefore anchored to the reduced form before any constraint is
switched on.

With the commitment constraints active, minimum stable output, ramp
limits, minimum up and down times and start-up costs stop the cheap CCGT from
following the residual perfectly, so the fast peakers pick up modestly more energy;
peaker energy rises by about 16~per cent, from 38.7 to 44.8~TWh across the
EU-27, the United Kingdom and Switzerland, and the fleet's full-load hours rise
from 139 on the reduced form's 277.7~GW fleet to about 160 on the unit commitment's own 280.5~GW fleet (Figure~\ref{fig:ucresults}a,b).

The extra running
cuts both ways. The hydrogen peaker burns more expensive fuel, but it also earns in
more hours, and on balance the missing-money position eases marginally. In the
\emph{unit-commitment case} the hydrogen peaker loses about \euro{}356.0~bn and the gas
peaker about \euro{}142.3~bn cumulatively to 2050, against \euro{}358.4 and \euro{}143.2
in the constraints-off limit case, so both positions improve, hydrogen by about \euro{}2.3~bn and gas by about
\euro{}0.9~bn, and the hydrogen-minus-gas gap narrows by
about \euro{}1.5~bn, from 215.2 to 213.7. The
two figures are intentionally distinct, the limit case being the constraints-off
reproduction of the reduced-form merit order and the unit-commitment case the constraints-on
result, and the small spread between them is the effect of the commitment constraints.

What the extension does not do is rescue the
capital-recovery verdict; at 160 hours the plant still recovers a small fraction of
its capital. The 5~per cent operating-reserve standard,
which the reduced-form merit order does not impose, implies a further about 59~GW of firm
capacity the system would need to add to meet the reserve margin
(Figure~\ref{fig:ucresults}c). This reserve adder is an adequacy signal distinct from,
and additional to, the 278~GW peak-capacity headline. Read against an installed
firm fleet of about 473~GW it is a roughly 13~per cent uplift, to an implied about
532~GW, the
kind of capacity-market or firm-build requirement a system planner would have to
procure beyond the energy-only peaker stack, and it leaves the 278~GW peaker headline
itself untouched. The 278~GW peak-capacity headline, the per-country win counts and
the qualitative verdict are all unchanged; the verdict now rests on a unit-commitment power model, and the hydrogen peaker still loses
heavily in every scenario.

\begin{figure}[H]
    \centering
    \caption{What the unit-commitment upgrade changes (EU-27+UK+CH, 2050). Panels~(a)
    and~(b) compare the unit-commitment run against the same run with its commitment
    limits switched off, both on that layer's own 280.5~GW fleet, so the left bar's 138
    full-load hours is not the 139 the reduced-form dispatch gives on its 277.7~GW fleet.
    Panel~(c) gives the per-country operating-reserve shortfall, whose 59~GW total is the
    firm capacity the reserve standard would have the system add.}
    \includegraphics[width=\textwidth]{paper/P34_uc_results.pdf}
    \label{fig:ucresults}
    \\[2pt]{\footnotesize\textit{Source: Authors' contribution.}}
\end{figure}

One mitigation deserves note. In combined-heat-and-power (CHP) mode the turbine's
recovered exhaust heat (raising total efficiency to approximately 85 per cent) can be
sold to district heating, and a winter peaker runs precisely when heat demand peaks.
Crediting that heat, scaled by each country's district-heating share and valued at
\euro{}50 to 100/MWh-heat, shrinks the loss only modestly (under H2~Push the gas
fleet moves to approximately \euro{}125 to 134~bn and the hydrogen fleet to approximately \euro{}342 to 350~bn), because a
\emph{power} peaker runs too few hours for the heat by-product to be large, and most
of Europe lacks the district-heating density to absorb it. The fuller cogeneration
case, a \emph{heat-led} unit running through the heating season, competes inside the
district-heating merit order, where hydrogen ranks last
(Appendix~\ref{sec:results:dhmerit}).

\begin{figure}[H]
    \centering
    \caption{Cumulative peaker profit to 2050, hydrogen against gas (EU-27+UK+CH).}
    \includegraphics[width=\textwidth]{paper/P22_peaker_cumulative.pdf}
    \label{fig:peakercum}
    \\[2pt]{\footnotesize\textit{Source: Authors' contribution.}}
\end{figure}


\subsection{Switzerland: Hydrogen Imports and Buildings Decarbonisation (RQ2)}
\label{sec:results:switzerland}

To address RQ2 we integrate Switzerland into the analysis \emph{endogenously}, closing
the three couplings that a one-directional exogenous check leaves open. First,
the Swiss delivered-hydrogen cost is made consistent with the EU supply side. It is the
bottom-up EU supply cost plus a pipeline-transport adder for the import distance (a
representative 600~km from the major European hydrogen producers), never read from a
standalone OIES import-price trajectory, so the Swiss molecule price and the EU ensemble
are mutually consistent. Second, a fixed-point feedback couples incremental Swiss
heat-pump load to the EU winter electricity price and back. Extra Swiss demand raises
the EU price, which raises the Swiss heat-pump cost, which feeds back to demand.

Because Switzerland is only about 2.27~per cent of the coupled demand and is
inter\-connector-bound, this loop converges to a tolerance of \euro{}0.01/MWh in a single
price update (the winter price moves once, from \euro{}240 to \euro{}246/MWh-e in 2050,
and is then stationary; the iteration history is reported in the Swiss integration
outputs in the supplementary material). Third, a seasonal-hydro credit recognises the Swiss
reservoir fleet (about 8~GW of firm winter headroom) as cheap firm winter capacity.
Together these make the Swiss case a member of the EU ensemble in place of a standalone
sensitivity, and the result generalises the finding to a high-renewable, cavern-poor,
import-reliant grid with its own seasonal hydro flexibility.

\begin{figure}[H]
    \centering
    \caption{Switzerland: hydrogen-boiler levelised cost against the heat-pump band. The three hydrogen lines are the stranded, central and rapid delivered-price paths, which Table~\ref{tab:scenarios} assigns to the four scenarios. The band spans the air-source and ground-source heat
    pump, and the dotted line is the gas boiler. Every series is Swiss, on the uncoupled
    arm that takes EU prices without feeding Swiss load back.}
    \includegraphics[width=0.72\textwidth]{paper/P14_switzerland.pdf}
    \label{fig:fig14}
    \\[2pt]{\footnotesize\textit{Source: Authors' contribution.}}
\end{figure}

The Swiss case runs one way on both margins, and the reason is the same on each. On
base load the heat pump wins comfortably against the incumbent fossil boiler and
narrowly against hydrogen by 2050; at the winter peak hydrogen loses in every scenario
once the Swiss reservoir fleet is credited as firm winter capacity, and wins under
H2~Push on the uncoupled 29-country series that does not credit it.

The Swiss
grid is predominantly hydro and nuclear, at 50~gCO$_2$/kWh in 2025 and on a
trajectory to near-zero by 2050, and against gas the heat pump is already the clear
2025 choice, at \euro{}160.5/MWh (air-source) and \euro{}175.8/MWh (ground-source)
against a Swiss gas-boiler cost of \euro{}216.0/MWh.

The ordering within the
heat-pump category matters here as it does EU-wide, since the air-source unit, and not the ground-source one, is the cheaper of the two once equipment capital is charged against peak sizing, which annual output does not enter. By 2050 the air-source unit falls to \euro{}122.4/MWh and the ground-source unit to
\euro{}131.1/MWh, and the gas boiler to \euro{}165.9/MWh, so the fossil comparison is
not close. The base-load band quoted here and in the next paragraph is the uncoupled
arm. We report it rather than the endogenous one because it is the narrower of the two
and so the conservative choice for the conclusion the section draws. Coupling the Swiss
load to the EU winter price moves the 2050 air-source unit to \euro{}120.9/MWh and the
hydrogen boiler to \euro{}130.6/MWh, widening the heat pump's margin from \euro{}5.6 to
\euro{}9.6/MWh. The peaking figures below are the endogenous arm's own, since that is
the test the coupling was built for.

The hydrogen comparison is closer, and it turns on Swiss geology. On the
EU-supply-consistent delivered cost, with the 600~km pipeline adder in place, the
hydrogen-boiler LCOH in Switzerland falls from \euro{}298.3/MWh in 2025 to
\euro{}128.0/MWh by 2050, which leaves the best Swiss heat pump \euro{}5.6/MWh the
cheaper (Figure~\ref{fig:fig14}). Switzerland has no domestic salt caverns, so a
Swiss hydrogen heating load is held across the season at about \euro{}3.0/kg rather
than the \euro{}1.5/kg the cavern markets pay, and that charge is what holds the
ordering. The opposite ordering, with hydrogen crossing below the
Swiss heat-pump band under the central supply trajectory, is an
artefact of omitting seasonal storage from the base-load comparison while charging
it in the peak and dispatch tests, so that the same molecule was priced two
different ways, and it does not survive the correction.

Closing the price feedback moves these levels only slightly, and in hydrogen's
disfavour. The fixed point raises the
Swiss winter electricity price by about \euro{}6/MWh-e, from \euro{}240 to
\euro{}246/MWh-e in 2050, and crediting the reservoir fleet as firm winter capacity
brings it down to about \euro{}141/MWh-e. The net effect on the levelised cost of
heat is between \euro{}1 and 3/MWh. The heat-pump LCOH is about \euro{}1.5/MWh
\emph{cheaper} once the seasonal-hydro credit and the consistent EU price are
applied, while the hydrogen route is about \euro{}2.5/MWh \emph{dearer} once its
cost is tied to the EU supply side in place of the more optimistic standalone
trajectory. Both movements widen the Swiss heat pump's margin, so the feedback works with the ordering. It is nonetheless of the same order
as the \euro{}5.6/MWh margin itself, which is why the Swiss base-load result is
reported here as the weaker of the two Swiss arguments and the peaking verdict is
left to carry the case.

What the endogenous treatment confirms most emphatically is the peaking verdict,
which is where the Swiss argument against hydrogen in heating has its margin. Cheap
firm Swiss hydro beats the imported hydrogen peaker in every one of the four
scenarios. The Swiss hydrogen peaker prices at \euro{}134.8 to 235.9/MWh-heat against a
heat-pump peak of \euro{}62.3 to 90.3/MWh-heat. That credit is what decides it. The
uncoupled 29-country series, which prices the Swiss peak at \euro{}155/MWh-heat without
the reservoir fleet, gives hydrogen the Swiss peak under H2~Push, so Switzerland sits
inside the count of 16 that series produces and outside the Swiss verdict this
subsection reports.
Switzerland's lack of domestic salt-cavern storage leaves the seasonal-storage
adder at about \euro{}100/MWh-heat, roughly double the cavern-endowed level, and
that adder is what shuts the peaking niche of
Section~\ref{sec:results:meritorder} in the Swiss winter peak in every scenario.
The base-load and peaking results now share that mechanism. The same absent geology that prices the Swiss peaker out of the
cold snap prices the Swiss seasonal store at twice the cavern rate on the base-load
comparison, where it is spread thinly over the whole year's heat without concentrating on the peak, which is why it decides the base-load case by
\euro{}5.6/MWh and the peaking case by an order of magnitude more.

The Swiss Energy Perspectives 2050+ \citep{SFOE2023perspectives} project heat pumps and
direct electrification as the dominant Swiss building-decarbonisation pathway, and the
model reaches that pathway on every margin it tests. Switzerland's geographic proximity
to major EU hydrogen producers provides a genuine cost advantage through shorter
pipeline distances \citep{OIES2024hydrogenimportsET32}, and on 2050 base-load cost that
advantage is nearly, but not quite, sufficient to overcome the efficiency advantage of
the heat pump on a near-zero-carbon grid. It comes as close as it does only because
Swiss retail electricity carries a full levy stack while hydrogen is priced delivered and
untaxed, which is the same fiscal asymmetry that
qualifies the EU-wide result and is set out in Appendix~\ref{sec:limitations}.

The endogenous integration therefore generalises the EU pattern to Switzerland, and the
pattern it generalises is the storage one, a cold, high-price, import-connected but
cavern-free system falling on the heat pump's side of both the base-load line and the
peak line. The two judgement-call inputs, the 600~km representative pipeline distance
and the 8~GW of firm reservoir headroom, are both conservative towards hydrogen, so
relaxing either would narrow the Swiss base-load margin and never widen it, and that
margin is already narrow; the peaking verdict, by contrast, has margin to spare against
both (Appendix~\ref{sec:limitations}).

Hydrogen's residual role in Switzerland lies outside the base-load comparison
altogether, and two segments remain unquantified here: the hard-to-electrify historic
building stock in city centres constrained by noise or space, and seasonal energy
storage integrated with Swiss hydropower. Quantifying these niches requires
city-level spatial modelling and sector-coupling analysis beyond the current scope.


\subsection{Technology Mix}
\label{sec:method:techshares}

Technology shares evolve by interpolation from an observed 2025 baseline to the
scenario-specific 2050 targets (Table~\ref{tab:scenarios}), and the interpolation has
three parts that we set out in full here, because the mix is described elsewhere in this
paper as scenario-imposed and that is only two-thirds true.

The first part is the path. Each share moves as
$s_{t} = s_{2025} + (s_{2050}-s_{2025})\,\tau^{p_{i}}$, where $\tau$ is the fraction of
the 2025 to 2050 window elapsed and the shape exponent is drawn per technology each
sample: $U(0.7,2.0)$ for heat pumps, $U(0.8,2.2)$ for district heat and $U(1.0,2.5)$ for
hydrogen, with fossil and resistance reusing the district-heat draw and biomass the
heat-pump draw. A draw can therefore front-load or back-load a transition without changing
its endpoints, except for hydrogen, whose exponent is bounded below at one and so can only
back-load. The second part is a cost-responsive reallocation among the non-fossil
options. Their combined mass is redistributed by a softmax on levelised cost,
$w_{i} \propto \exp(-c_{i}/T)$ with $T = 50$~\euro{}/MWh, so a cheaper technology takes a
larger share without taking all of it. The third part is the weight given to that
reallocation, which ramps linearly from zero in 2025 to $0.50$ by 2050. The realised
share is the weighted blend of the interpolated path and the cost-driven allocation.

A fourth part applies to hydrogen alone and belongs here because it bounds a headline.
After the blend, hydrogen is capped at its own drawn 2050 target, any excess is routed to
heat pumps up to their feasibility ceiling and then to district heat, and the mix is
renormalised. The cap is one-directional, so cost can push hydrogen below its drawn target
but never above it. It exists because hydrogen enters this model as an exogenous adoption setting for a
peaking share rather than a base-load winner, and without it the softmax and the fossil-ban
reallocation inflate hydrogen well past its design value. Uncapped, the softmax realises about
13~per cent under Stated Policies against a 5~per cent setting. The statement that hydrogen
reaches at most about 8.5~per cent of useful heat in any scenario we run is therefore a
property of this cap and not an independent result of the cost mechanism.

Three further consequences follow and are load-bearing elsewhere. The mix is not purely a
scenario input, since by 2050 half the non-fossil allocation is decided by cost. Nor is
it purely cost-driven, since the fossil share is governed by the policy-ambition lever
throughout and the cost-responsive weight is only a tenth at 2030, which is why the
one-at-a-time tornado of Appendix~\ref{sec:sensitivity:tornado} returns such small
numbers at that date. And the softmax temperature is a judgement. We chose it so that a
cheaper technology wins moderately rather than absolutely, and we did not estimate it
from market data, which is the largest unestimated parameter in the demand-side layer.
Replacing this heuristic with a discrete-choice allocation estimated on observed
switching is the natural extension.

The 2025 baseline
shares are EU averages drawn from EHPA heat-pump market data and Euroheat \& Power
district-heating statistics: heat pumps approximately 18\,\%, district heat
approximately 12\,\%, biomass approximately 10\,\%, and fossil fuels (gas and
oil) approximately 60\,\%. These EU-average shares are stated on a different
basis from the per-country useful-heat shares reported in
Section~\ref{sec:results:techshares} (heat pumps 8~per cent, gas 44~per cent,
oil 14~per cent). The figures here are EU-aggregate EHPA and Euroheat \& Power statistics; the results section reports the per-country useful-heat
distribution actually used in the model, so the two are not contradictory.
For each NUTS3 region and building type, the 2050
heat-pump target is modulated by a feasibility score:
\begin{equation}
    \mathrm{HP_{target}}(r,b) = s_{\mathrm{HP}}\,\bigl(0.4 + 0.6\,f_{\mathrm{HP}}(b)\bigr),
    \label{eq:hpfeas}
\end{equation}
where $s_{\mathrm{HP}}$ is the scenario heat-pump target and
$f_{\mathrm{HP}}(b) \in [0,1]$ is a building-type feasibility factor
(single-family 0.9, multi-family or high-rise 0.5, other 0.6); the intercept of
0.4 ensures that even a zero-feasibility region still reaches 40\,\% of the
scenario target. District heat is weighted by an analogous feasibility factor
(single-family 0.3, multi-family or high-rise 0.8), reflecting the economics of
network heat at different building densities.

\paragraph{Feasibility-scoring rationale.} The feasibility factors are anchored to the
central tendency of the heat-pump retrofit and district-heating literature, synthesised
across several sources, and they are flagged as a key modelling assumption in the
Limitations section. Table~\ref{tab:feasibility} sets the adopted central values against
the indicative feasibility ranges reported in that literature, by building type; the
values used here sit inside the cited ranges. The ordering encodes two well-documented
patterns. First, air-source heat pumps retrofit most readily into detached single-family
homes, which have the facade access, emitter sizing headroom and outdoor-unit space that
low-temperature operation needs, and least readily into dense apartment blocks, where
shared distribution, planning consent and outdoor-unit siting raise the barrier; this is
the gradient reported across the heat-pump retrofit and barriers literature
\citep{Gibb2023Coming,Corbett2023Pumping}, hence the descending single-family $>$ other
$>$ multi-family ordering of $f_{\mathrm{HP}}$.

Second, district heat is economic where linear heat density is high, that is in compact
multi-family stock and uneconomic in dispersed single-family stock \citep{Persson2019Heat},
which inverts the ordering for the district-heat factor. The non-zero intercept of 0.4
reflects that no region is structurally barred from heat pumps. Even the least
favourable stock can reach a substantial share through ground-source units, communal and
high-temperature heat pumps and new build. Because these coefficients shape the 2050
heat-pump and district-heat shares, they are treated as a structural assumption. The
technology-share targets are sampled in the Monte Carlo
(Appendix~\ref{sec:method:uncertainty}) and the feasibility design is revisited in the
Limitations section (Appendix~\ref{sec:limitations}). The qualitative conclusions of the
paper turn on the merit-order economics rather than on the precise feasibility split,
since hydrogen's heating share is derived from the peak dispatch
(Appendix~\ref{sec:method:merit}) and is not set by these factors.

\begin{table}[H]
\centering
\caption{Heat-pump and district-heat feasibility factors by building type.}
\label{tab:feasibility}
\begin{tabularx}{\textwidth}{lXXX}
\toprule
Building type & Adopted $f_{\mathrm{HP}}$ & Indicative literature range
             & Basis \\
\midrule
Single-family (SFH)       & 0.9 & 0.7 to 0.95 & \citet{Gibb2023Coming,IRENA2022heatpumps} \\
Other / mixed             & 0.6 & 0.4 to 0.7  & \citet{Corbett2023Pumping,JRC2023smallheating} \\
Multi-family / high-rise  & 0.5 & 0.3 to 0.6  & \citet{JRC2023smallheating,IRENA2022heatpumps} \\
\midrule
 & Adopted $f_{\mathrm{DH}}$ & Indicative literature range & Basis \\
\midrule
Single-family (SFH)       & 0.3 & 0.1 to 0.4  & \citet{Persson2019Heat} \\
Multi-family / high-rise  & 0.8 & 0.6 to 0.9  & \citet{Persson2019Heat} \\
\bottomrule
\end{tabularx}
\smallskip
\raggedright\small
The ranges are indicative bands synthesised from the cited heat-pump retrofit,
adoption-barrier and linear-heat-density literature; they bracket the structural feasibility of each technology in each
building type. The heat-pump gradient (single-family $>$ other $>$ multi-family)
follows the retrofit and barriers evidence
\citep{Gibb2023Coming,Corbett2023Pumping};
the district-heat gradient inverts it because network heat is economic only at
high linear heat density \citep{Persson2019Heat}. Because the central values are
structural, the technology-share targets they shape are sampled in the Monte
Carlo (Appendix~\ref{sec:method:uncertainty}) and the design is revisited in the
Limitations section (Appendix~\ref{sec:limitations}).
\\[2pt]{\footnotesize\textit{Source: Authors' contribution.}}
\end{table}


\subsection{The Emissions Layer}
\label{sec:method:emissions}

The emissions inventory runs after the technology split, on the same
country-by-year final-energy quantities the cost layer prices. Each Monte Carlo draw
carries its own inventory, so the fan of Section~\ref{sec:results:emissions} is the spread of the sampled
technology mix, with no separate emissions uncertainty laid on top of it.
The scope is direct, scope-1 combustion at the building plus the indirect content of the
two delivered carriers, and it excludes embodied emissions in equipment, network
construction and hydrogen production upstream of the delivered price.

Combustion fuels carry fixed factors on a lower-heating-value basis of final energy,
natural gas at $0.202$ and heating oil at $0.265$~tCO$_2$/MWh, taken from Annex~II of the Sixth Assessment Report's Working Group~III volume and cross-checked against the UK government conversion factors. Biomass and hydrogen are entered at zero. Zero is the European accounting convention for a renewable fuel under the Renewable
Energy Directive; it is not a physical claim about the molecule, and at 17~per~cent of the 2025 mix the biomass
convention is material to the level of every pathway. It flatters all four equally,
because the biomass share is a scenario input common to them, so it moves the levels
without moving the ranking or the gap between Current~Policies and Net~Zero that
Section~\ref{sec:results:emissions} reports.

Electricity and district heat carry country-by-year intensities, because the whole
electrification case turns on the grid cleaning up. Both start from
observed 2025 country values \citep{EMBER2024electricity} and follow declining paths to
2050 that each scenario scales by its grid-decarbonisation lever. The electricity paths
run, for example, from 600 to 30~gCO$_2$/kWh in Poland, 350 to 15 in Germany, 80 to 8 in
France and 50 to 5 in Switzerland; district heat follows a parallel and generally lower
country path \citep{EEA2023fit55}. Values between the milestone years are linearly interpolated. Every country in the study
area carries its own series, so the implementation's fallback to the German path is
unreachable here.

Two properties of that construction bound what the emissions results can carry, and both
are set out again with their direction in Appendix~\ref{sec:limitations}. The post-2025
intensity paths are scenario-consistent trajectories calibrated to published 2025 levels
and to European decarbonisation targets, not published country projections, so the
absolute 2050 levels inherit that calibration while the ranking across scenarios does not,
since every scenario reads the same underlying path through its own lever. And the
intensity is an annual average applied to all heating electricity, while the cost layer
prices the heat pump's winter draw at an hourly scarcity price. Heat-pump peak electricity
is both dearer and dirtier than the annual mean, so the inventory carries a systematic
bias in the heat pump's favour that no hourly ledger in this model can quantify.

\subsection{Techno-Economic Assumptions}
\label{sec:method:techassumptions}

Table~\ref{tab:techparams} summarises the key techno-economic parameters for each
technology. All monetary values are in EUR$_{2024}$.

\begin{table}[H]
\centering
\caption{Techno-economic parameters by heating technology (EUR$_{2024}$).}
\label{tab:techparams}
\begin{tabularx}{\textwidth}{lXXXX}
\toprule
Technology & CAPEX 2025 & CAPEX 2050 & Efficiency / SCOP & Fuel \\
           & (EUR/kW)   & (EUR/kW)   & 2025 / 2050      &      \\
\midrule
Gas boiler         & 420     & 400     & 0.92 / 0.94 & Natural gas \\
Oil boiler         & 450     & 430     & 0.90 / 0.90 & Oil \\
Biomass boiler     & 950     & 850     & 0.88 / 0.90 & Biomass \\
Resistance heater  & 200     & 180     & 0.98 / 0.99 & Electricity \\
HP (air-source)    & 1{,}200 & 800     & 3.3 / 3.5   & Electricity \\
HP (ground-source) & 2{,}000 & 1{,}400 & 3.8 / 4.3   & Electricity \\
District heat      & 500     & 450     & 0.95 / 0.97 & District heat \\
H$_2$ boiler       & 600     & 550     & 0.90 / 0.92 & Hydrogen \\
\bottomrule
\end{tabularx}
\smallskip
\raggedright\small
CAPEX and efficiency values are anchored to the Danish Energy Agency
Technology Catalogue installed-cost figures \citep{DEA2023technology},
cross-checked against the JRC Technology Data repository
\citep{JRC2023technology,JRC2023smallheating}, the IEA \emph{Future of Heat
Pumps} \citep{IEA2022heatpumps} and IRENA \citep{IRENA2022heatpumpcosts}, through the
source audit set out in Appendix~\ref{sec:costopt}. The Danish Energy Agency
catalogue is the primary anchor for every technology because it reports installed (turnkey) costs in place of equipment-only costs on a consistent basis; the
other sources are used as cross-checks and to set the heat-pump and hydrogen
values where the catalogue is less granular. The two tabulated years are the
endpoints of the cost path. Intermediate-year CAPEX is obtained by linear
interpolation between the 2025 and 2050 values, so the table fully specifies the
trajectory, and no separate learning-curve exponent is applied. Heat-pump
performance is the seasonal coefficient of performance (SCOP); the values shown
are EU averages following EHPA data \citep{EHPA2023scop} and EU Regulation
813/2013, with the air-source SCOP set to the EHPA/Eurovent field value (3.3).
The single SCOP figure is an EU-average input only. Because the seasonal
coefficient of performance falls in colder climates and at higher flow
temperatures, the model applies a temperature-dependent derate at the country
level through the cold-snap performance $\mathrm{COP}^{\mathrm{cold}}_{c}$ and
the hourly dispatch (Appendix~\ref{sec:method:merit}). CAPEX is additionally
scaled by a country-specific construction labour-cost multiplier derived from
Eurostat \texttt{lc\_lci\_lev} (2024), and is country- and year-resolved in the
model.
\\[2pt]{\footnotesize\textit{Source: Authors' contribution.}}
\end{table}


\subsection{Fuel Price Assumptions}
\label{sec:method:prices}

Residential end-user energy prices for the 2025 baseline are taken from Eurostat:
natural gas from dataset \texttt{nrg\_pc\_202} and electricity from
\texttt{nrg\_pc\_204}, both for the first half of 2025. Heating-oil, biomass and
district-heat prices are per country as well, assembled from the EC Weekly Oil Bulletin
and GlobalPetrolPrices, from ENplus, proPellets and national pellet series, and from
Euroheat and national heat regulators; their provenance is lighter than the two Eurostat
carriers and we treat it as a limitation. The green-hydrogen price follows the European
Hydrogen Observatory (2024). Table~\ref{tab:prices} gives the 2025 levels averaged over
the 29 markets and the multipliers applied to reach 2050. Prices vary substantially
between countries. Residential gas, for example, ranges from about \euro{}31/MWh in
Hungary to \euro{}213/MWh in Sweden.

\begin{table}[H]
\centering
\caption{Energy-carrier prices: 2025 baseline and 2050 multiplier.}
\label{tab:prices}
\begin{tabularx}{\textwidth}{lXXX}
\toprule
Carrier & 2025, mean of the 29 markets (EUR/MWh) & 2050 multiplier & Source \\
\midrule
Natural gas   & 103 & 0.55      & Eurostat \texttt{nrg\_pc\_202} (H1 2025) \\
Electricity   & 257 & 0.80      & Eurostat \texttt{nrg\_pc\_204} (H1 2025) \\
Heating oil   & 154 & 0.55      & EC Weekly Oil Bulletin, GlobalPetrolPrices \\
Biomass       & 69  & 1.05      & ENplus, proPellets, national series \\
District heat & 88  & 0.80      & Euroheat, national regulators \\
Hydrogen      & 200 & see text  & European Hydrogen Observatory (2024) \\
\bottomrule
\end{tabularx}
\smallskip
\raggedright\small
Every carrier is priced per country, and the column above is the unweighted mean of
those 29 values, so it reconciles with the per-country prices the model runs on. It is
not the Eurostat headline aggregate, which is weighted by consumption and reads
\euro{}114.3/MWh for gas and \euro{}287.2 for electricity; those are the figures the
2025 parity discussion of Section~\ref{sec:results:lcoh} quotes, and the difference is the weighting; the vintage is the same. The hydrogen price follows scenario-specific trajectories. The central trajectory falls from
\euro{}200/MWh in 2025 to \euro{}90/MWh in 2030 and \euro{}50/MWh by 2050. For
Switzerland, hydrogen import prices are treated as a separate key sensitivity variable
(Appendix~\ref{sec:method:uncertainty}). The oil, biomass and district-heat series carry
lighter provenance than the two Eurostat carriers, and five district-heat entries (Spain,
Portugal, Ireland, Cyprus, Malta) are neutral placeholders for markets whose
district-heat share is near zero.
\\[2pt]{\footnotesize\textit{Source: Authors' contribution.}}
\end{table}


\subsection{Hydrogen Supply Cost and the Last-Mile Delivery Networks}
\label{sec:method:h2}

Hydrogen requires two country-specific treatments beyond the hub price trajectory:
where the molecule comes from (the delivered commodity cost) and how it reaches the
dwelling (the distribution network). Electric heating requires the symmetric
treatment of the second, since its winter peak also drives incremental last-mile
capital, and the two adders defined below are both active in the headline levelised
costs so that each carrier is charged for its own delivery network.

\paragraph{Per-country delivered-cost multiplier.} The hub trajectory is scaled to each
country by a delivered-cost multiplier $m_c$, so the residential hydrogen price is
$p_{c,t} = p^{\mathrm{hub}}_t \cdot m_c$. Rather than a uniform or heuristic assumption,
$m_c$ is derived from a bottom-up delivered-cost model that, for each country, levelises
every supply route, namely domestic \emph{green} electrolysis, \emph{pink} (nuclear)
electrolysis, \emph{blue} (gas reforming with CCS), pipeline \emph{import}, and seaborne
\emph{ammonia} import, and takes the cheapest. Each route depends on the country's cost
of capital (WACC), construction labour cost, electricity or feedstock price, renewable
capacity factor, water scarcity, pipeline connectivity, and nuclear or
CO\textsubscript{2}-storage availability, so the multiplier captures cost drivers beyond
renewable resource quality alone.

The technology constants are calibrated to verified 2050 benchmarks: electrolyser CAPEX
(approximately \euro{}300/kW) and green LCOH from IRENA \citep{IRENA2022trade}; blue
(approximately \euro{}2.7/kg) from IEAGHG \citep{IEAGHG2022blue}; pipeline transport
(approximately \euro{}0.16/kg per 1{,}000~km) from the European Hydrogen Backbone
\citep{EHB2023countries}; and delivered or import costs from OIES
\citep{Alsulaiman2023ET24,OIES2024hydrogenimportsET32}. The multiplier is the country's
delivered cost divided by the EU-median delivered cost, so the median country pays the
hub price ($m_c=1.0$). The resulting $m_c$ ranges from 0.77 (Ireland: strong wind, low
capital cost) to 1.33 (Cyprus: an isolated island with no pipeline); pink and blue are
not the cheapest route in 2050. Each multiplier is held as a low/central/high band,
with the sensitivity of the central findings reported in
Appendix~\ref{sec:sensitivity:h2supply} and the full derivation and sources in the
supplementary material.

\paragraph{Hydrogen distribution cost.} Because a methane network cannot carry
100\% hydrogen unmodified, the hydrogen-boiler LCOH carries a distribution adder
$a_c$ (EUR/MWh\textsubscript{useful})
\citep{NorthernGasNetworks2018h21,GasForClimate2020backbone}. A per-dwelling
capital cost $K_c$ is annualised over a 40-year network life at the country
discount rate $r_c$ and spread over the annual useful heat per dwelling
$Q\approx13$~MWh:
\begin{equation}
a_c = \frac{\mathrm{CRF}(r_c,\,40)\cdot K_c}{Q},
\label{eq:h2adder}
\end{equation}
where $\mathrm{CRF}$ is the capital-recovery factor. The per-connection cost is
genuinely uncertain, so $K_c$ is carried as a low/central/high range (conversion
\euro{}1{,}000/1{,}500/3{,}000 per dwelling; new-build \euro{}3{,}000/5{,}000/8{,}000) and
across three network modes, namely \emph{retrofit} conversion of the existing grid,
\emph{new-build} of a dedicated network, and the coverage-weighted \emph{blend} of
the two. The adder is \emph{on by default}, in blend mode and at the central cost
bound, so every headline levelised cost reported in this paper already carries it;
free reuse, retrofit and new-build are reported as a sweep around that default
(Appendix~\ref{sec:results:h2infra}).

\paragraph{Seasonal storage on the base-load hydrogen load.}
\label{sec:method:storageadder} A hydrogen heating load has
to be held from the season in which the molecule is cheapest to make to the season in
which it is burned, and the base-load levelised cost charges it for doing so. On the
model's own monthly shape, space heating draws a median 0.748 of its annual energy in
the October-to-March half of the year, within a range of 0.710 to 0.787 across the 29
countries, whereas electrolytic production is far flatter, so the difference between the
two profiles is the part of the load that must be time-shifted.

We charge the seasonal store on that difference, at each country's own winter share less
one half, a fraction whose median across the study area is 0.25. The store itself is
priced at the same rates the peak and dispatch tests use, \euro{}1.5/kg where domestic
salt caverns are available and \euro{}3.0/kg where they are not, at about one cycle a
year \citep{Caglayan2020caverns}. The base-load charge divides those rates by the
year-specific boiler efficiency, 0.92 at 2050, where the dispatch tests divide by the
0.90 of their own peaking convention, so per unit of useful heat this is worth
\euro{}12.1/MWh in cavern-endowed Germany and \euro{}28.4/MWh in cavern-free Spain, and it is that
difference, and not the delivered price of the molecule, that separates the countries
in which hydrogen holds base load from those in which it does not.

An asymmetric accounting charges this term in the building-peak and power-dispatch tests
but not in the base-load levelised cost, so the same molecule is priced two different ways
and the base-load comparison omits a cost the dispatch comparison imposes. We charge it on
both, and the choice of fraction is deliberately conservative towards
hydrogen in two respects. It credits hydrogen with meeting the whole non-seasonal half
of its load from contemporaneous production, which no real supply chain would achieve
exactly, and it ignores the within-winter swing, which also has to be stored somewhere.
A fuller accounting would only charge more. The sensitivity is correspondingly
one-sided, and because the count is close to linear in the charged fraction we sweep it
continuously. Charging nothing leaves the best heat pump cheaper
in 12 of 29 countries at a median gap of $-$\euro{}5.1/MWh.

A quarter of the load-derived fraction gives 15 countries at \euro{}0.0/MWh, half gives
19 at $+$\euro{}5.2, three-quarters gives 21 at $+$\euro{}10.3, and the fraction itself,
which is the basis adopted here, gives 22 at $+$\euro{}15.4. Above the adopted basis the
count keeps climbing, to 25 at one and a half times and 26 at twice, and charging about
four times the load-derived fraction leaves the heat pump cheaper in all 29 at $+$\euro{}81.1/MWh, with hydrogen winning nowhere. That last arm is a multiple of each
country's own fraction and not a flat full-storage charge, so it bills the southern
markets slightly more than their whole load and the Nordics slightly less; charging the
full adder everywhere instead gives $+$\euro{}87.8/MWh and the same count of 29. The adopted basis therefore
sits near the low end of the range and well below what a fuller accounting would charge,
and the defensible summary of the base-load verdict is 22 countries on that basis, 15 to
26 across the fractions we regard as plausible.

\paragraph{Electricity distribution reinforcement.} Electric heating carries the
symmetric charge. Mass heat-pump deployment raises the coincident winter peak that
the low- and medium-voltage network must carry, and we treat that incremental
reinforcement as not already recovered in today's retail tariff. We therefore add a
reinforcement adder
$a^{\mathrm{el}}_{c}$ (EUR/MWh\textsubscript{useful}) to the levelised cost of every
electric heating technology. The added peak per dwelling
$P^{\mathrm{dw}}_{c}$ is the annual useful heat divided by the country's
climate-scaled full-load hours $h_c$ and by the cold-snap coefficient of performance
a cold-snap coefficient of performance $\mathrm{COP}^{\mathrm{el}}=2.3$, and it is multiplied by
a coincidence factor $\kappa=0.5$ so that dwellings are not assumed to draw their
maxima at the same instant:
\begin{equation}
a^{\mathrm{el}}_{c}
    = \frac{\mathrm{CRF}(r_c,\,40)\cdot u^{\mathrm{el}}\cdot
            \kappa\, P^{\mathrm{dw}}_{c}}{Q},
\qquad
P^{\mathrm{dw}}_{c} = \frac{Q}{h_c\,\mathrm{COP}^{\mathrm{el}}},
\label{eq:eladder}
\end{equation}
with $u^{\mathrm{el}}\in\{400,800,1{,}600\}$~EUR per added kW of peak (low, central
and high), and with the same 40-year network life and country discount rate as the
hydrogen adder, so the two carriers are annualised on identical terms. Across the
band the adder is worth \euro{}1.1 to 38.2/MWh, and it is highest in peaky, mild-climate
systems where a small annual heat demand is delivered through a short, sharp
season. Like the hydrogen adder it is \emph{on by default} at the central bound.
Charging each carrier its own last-mile network removes an asymmetry, in which
hydrogen would pay for its pipes only inside a dedicated scenario while electric
heating never paid for its wires. It also introduces a
residual risk, that retail electricity tariffs already recover part of today's
network cost so that the incremental charge may double-count, which is stated in
Appendix~\ref{sec:limitations}.


\subsection{District Heating: Source-Switching Merit Order}
\label{sec:results:dhmerit}

The peaking test extends to a third scale, the district-heating peak, where the same
merit-order logic governs hydrogen inside a single operator's stack of sources
dispatched into the same pipes. Ranked by marginal cost,
the 2050 stack runs: waste and excess heat (must-run, approximately \euro{}10/MWh),
large heat pumps, biomass CHP, gas CHP (carbon-priced), and hydrogen. Hydrogen's
rank in the stack depends on its delivered price. With the expensive hydrogen of the
low-carbon scenarios it is the dearest source and catches only the residual peak; under the cheap derived hydrogen of H2 Push it ranks just below gas CHP. It
undercuts gas CHP on operating cost in 0, 0, 7 and 7 of 29 countries across the
four scenarios on the symmetric basis, and in 0, 5, 11 and 20 on the residential-tariff
basis the next paragraph sets out; gas CHP nonetheless remains the marginal source that
sets the district-heat peak in 20 of 29 (hydrogen in 9), since hydrogen is dispatched
below it, so 20 is the arena's structural ceiling on either basis.
As at the building level, it is the carbon price on the gas alternative,
reinforced by cheap hydrogen, that opens the door, and the addressable share is the
peak sliver above the waste/heat-pump/biomass base load, an EU-weighted 0 to 12 per cent of district heat (Figure~\ref{fig:dhstack}).

One price basis carries this arena and should be stated. The stack charges gas at the
Eurostat residential retail tariff including taxes, a study-area median of
\euro{}55/MWh in 2050, because that is the series the whole cost layer runs on; a
district-heat operator buys wholesale. Repricing the gas unit at the \euro{}30/MWh hub
price the power block already uses takes the four counts from 0, 5, 11 and 20 to 0, 0, 7
and 7. The test compares hydrogen against gas alone, so the electricity price the large
heat pump pays does not enter it. We headline the 7, because it is what a wholesale-buying
operator would see and it is the figure a district-heat investor should use, and report
the 20 alongside it as what the residential basis returns.

\begin{figure}[H]
    \centering
    \caption{District heating. Panel (a) is the merit-order stack under H2~Push, the
    scenario most favourable to hydrogen; panel (b) is the Net Zero build-out. The two
    panels are on different scenarios deliberately, so the stack is read at hydrogen's
    best and the expansion at the deepest decarbonisation.}
    \includegraphics[width=\textwidth]{paper/P12_dh_stack_expansion.pdf}
    \label{fig:dhstack}
    \\[2pt]{\footnotesize\textit{Source: Authors' contribution.}}
\end{figure}

\paragraph{Hourly heat dispatch.} Resolving the same stack to the hour
(Figure~\ref{fig:heatdispatch}) confirms the picture and adds the heat-led case. On the
cold-snap day, the decarbonised base of waste heat, large heat pumps and biomass carries
the bulk, and hydrogen fills the morning and evening peaks, the heat-side mirror of the
electricity peaker. This is hydrogen's most favourable framing; as a \emph{heat-led}
cogeneration unit it can run many winter hours, against the about 139 of the power
peaker, so its capital is far better utilised.

Even so, in the H2-favourable scenario it reaches only 9.3 per cent of EU district-heat energy as an aggregate, hydrogen's total over the study area's total (14.3 per cent on the single representative cold-snap day of Figure~\ref{fig:heatdispatch}, which loads the high-share markets), approximately 29 per cent in the three salt-cavern markets where
its delivered-plus-storage cost undercuts gas CHP (Germany, the United Kingdom and
Denmark), but a marginal 0.33 per cent across the other 26 on a demand-weighted basis, and at most 0.5 per cent in any single one of them, including cavern-endowed Poland, where cheaper sources fill the district-heat peak. The dispatch runs on the Stated Policies district-heat footprint, in which district heating is 18.6 per cent of residential heat, so this best case is 1.7 per cent of total EU heat. Both figures are taken on that one footprint. The H2~Push district-heat share is higher, at 22.9 per cent, and combining it with a dispatch solved on the Stated Policies footprint would mix two scenarios. Cogeneration is therefore hydrogen's strongest
niche in heating, but a narrow one, and one realised only where cheap hydrogen meets
salt-cavern storage.

\begin{figure}[H]
    \centering
    \caption{Hourly district-heat dispatch on the cold-snap day (EU-27+UK+CH, 2050). The
    marginal costs that order the stack are H2~Push; the district-heat footprint the stack is scaled to is Stated~Policies, so the two are on different scenarios and the shares in the paragraph above are read on the second. Waste and excess heat is a constant band because it is a
    by-product flow that does not follow the load. Gas cogeneration appears in seven hours,
    four across the morning peak and three across the evening one, and it sets the price in
    each of them.}
    \includegraphics[width=\textwidth]{paper/P23_dh_heat_dispatch.pdf}
    \label{fig:heatdispatch}
    \\[2pt]{\footnotesize\textit{Source: Authors' contribution.}}
\end{figure}

%── A.15 The infrastructure bill ───────────────────────────────────────────────

\subsection{The Network-Infrastructure Bill of Each Pathway}
\label{sec:results:infrabill}

Having tested hydrogen on the base-load, peaking and investment margins, we close
with the cross-pathway comparison those per-technology tests omit, namely the
network capital each whole pathway commits. Each pathway carries a different
wires-and-pipes bill that per-MWh cost comparisons obscure. We aggregate three
network capital
requirements from the model's own technology shifts: electricity distribution
reinforcement (the added coincident
winter heat-pump peak, at \euro{}400/800/1{,}600 per added kW of peak),
district-heat network expansion (new connections at \euro{}3{,}000/5{,}000/8{,}000
per dwelling), and the hydrogen distribution network (the coverage-weighted
conversion and new-build blend of Appendix~\ref{sec:results:h2infra}).

\begin{figure}[H]
    \centering
    \caption{The cumulative 2025 to 2050 network-infrastructure bill of each pathway, all 29
    markets. The number beside each bar is the central estimate; the whiskers are the low
    and high unit-cost cases. Every pathway is dominated by district-heat expansion and
    grid reinforcement, which a hydrogen network does not drive.}
    \includegraphics[width=0.7\textwidth]{paper/P13_infrastructure_bill.pdf}
    \label{fig:infrabill}
    \\[2pt]{\footnotesize\textit{Source: Authors' contribution.}}
\end{figure}

The cumulative 2025 to 2050 bills (central [low to high]), costed across all 29
markets, are \euro{}249~bn [143 to 423] under Current Policies, \euro{}325~bn [187 to 558] under Stated Policies, \euro{}506~bn [292 to 861] under Net Zero, and \euro{}412~bn
[240 to 705] under H2 Push. Two findings stand out. First, \emph{the
deep pathway's bill is dominated by district heating}, at \euro{}356~bn of Net Zero's
\euro{}506~bn against \euro{}119~bn of electricity reinforcement, so deep
decarbonisation is a district-heating build-out as much as a heat-pump roll-out.
Second, \emph{hydrogen does not buy its way out of grid costs}. H2 Push pays
\euro{}74~bn for a hydrogen network and still requires \euro{}89~bn of electricity
reinforcement, because
hydrogen serves only the peak sliver while heat pumps carry the bulk of the load,
so the hydrogen network bill is additive to the electrification bill and substitutes for no part of it. The same pattern holds in reverse across the sweep, since even
Current Policies, with the smallest hydrogen network at \euro{}4bn, still carries
\euro{}62~bn of reinforcement and \euro{}183~bn of district-heat expansion.

A complementary robustness check charges the heat pump the \emph{high} reinforcement
bound in the levelised-cost comparison while charging hydrogen nothing at all for
its network, which is the most conservative treatment of electric heating we can
construct. Under it heat pumps are cheaper than hydrogen in 10 of 29 countries and
cheaper than gas in 21 of 29, and the mean heat-pump-to-hydrogen headroom is
$-$\euro{}6.2/MWh, meaning that hydrogen leads on average by that margin
(Appendix~\ref{sec:sensitivity}). Restoring hydrogen's own network charge, so that
both carriers pay, returns the count to 22 of 29. The base-load ranking therefore depends on how the two last-mile networks are charged,
and the robustness statement is weaker for it.

